\documentclass[10pt,a4paper]{article}

\usepackage[T1]{fontenc}
\usepackage{lmodern}
\usepackage{microtype}
\usepackage[a4paper, margin=1.05cm]{geometry}
\usepackage{xcolor}
\usepackage{hyperref}

\definecolor{links}{RGB}{15, 76, 129}
\definecolor{primary}{RGB}{0, 0, 0}

\hypersetup{
    colorlinks=true,
    linkcolor=links,
    urlcolor=links,
    citecolor=links
}

\setlength{\parindent}{0pt}
\setlength{\parskip}{3pt}
\setlength{\leftmargini}{12.7pt}

\begin{document}

\begin{center}
    {\Large \textbf{\textcolor{primary}{The Budget Illusion: Why Bollywood's ``Safe Bets'' May Not Be Safe}}}\\[2pt]
    {\small Open Project $|$ CS215: Data Analysis and Interpretation $|$ IIT Bombay}\\[2pt]
    {\small \textbf{Checkpoint 1 Progress Report}}\\[1.5pt]
    \rule{\textwidth}{0.4pt}
\end{center}
\vspace{-10.5pt}

\textbf{\textcolor{primary}{Section 1: Meta Information}}

\vspace{1pt}
\textbf{Group 05:} \hfill \textbf{Mentor:} Avadh Hindocha

\textbf{Members:}\\
Ajay -- 25B0929\\
Bhargav -- 25B0940\\
Rushmith -- 25B1029\\
Aryan -- 25B1098

\vspace{12pt}
\textbf{GitHub Repository:} \url{https://github.com/Chaos0Speed/IndianMovies/}

\vspace{5pt}
\textbf{\textcolor{primary}{Section 2: Progress So Far}}

\vspace{1pt}
\textbf{Descriptive Statistics of the Dataset}\\
We use the Bollywood Movies dataset by Prashant Premkumar (2020), obtained from Mendeley Data (\href{https://data.mendeley.com/datasets/3c57btcxy9/1}{dataset link}). The
primary dataset consists of 1,698 rows and 14 columns detailing various attributes of Hindi films. Key columns include
budget, revenue (in INR), screen count, genre, release period, remake status, franchise status, and cast/crew novelty. The
dataset was obtained from public domain sources (publicly accessible data permissions) ensuring no access or compliance
issues. Initial inspection revealed the absence of release years, which restricts our ability to track temporal inflation.

\textbf{Notebook 1 -- Data Cleaning}\\
We conducted rigorous data cleaning to establish a reliable foundation. Duplicate rows were removed, and string inconsistencies
across categorical values were standardized. We converted categorical Yes/No variables (such as remake, franchise, new actor,
new director, and new music director) and holiday/normal release periods into binary indicators. Genres were split and
one-hot encoded using Pandas dummy variables. To account for film scale, we engineered a composite `Movie\_Scale` feature
by log-transforming budget and screen counts, normalizing them, and combining them with equal weights. Additionally, we
calculated two critical success metrics: revenue per screen and the revenue-to-budget ratio. The final processed data was
exported as `CleanedDataSet.csv` for downstream analysis.

\textbf{Notebook 2 -- Exploratory Data Analysis}\\
Our exploratory data analysis involved examining data types, checking unique categories, and compiling numerical summaries
for our engineered dataset. We visualized category counts and percentages to map the distribution of holiday releases,
remakes, franchises, and debutant talent. Crucially, we plotted histograms to observe the distributions of numerical features
like budget, revenue, and revenue-to-budget ratios. The histograms clearly show significant right-skewness in financial returns,
with a heavy tail of blockbusters masking a large cluster of average/low-performing movies. These histogram visualizations
underscore the necessity of applying log transformations or non-parametric tests when analyzing success metrics across
categories (e.g., new versus established talent).

\textbf{Fine-grained Problem Statement Refinement}\\
We aim to investigate whether traditional Bollywood “safe bets” specifically holiday releases, remakes, and franchise
continuation genuinely yield better financial performance (measured by revenue-to-budget ratio and revenue per screen).
Furthermore, we are evaluating whether introducing debutant talent impacts a movie’s financial success once we control for
budget scale and screen count. We will use group comparisons, hypothesis testing, and regression analysis to determine if
these factors remain significant predictors of success beyond raw financial backing.

\textbf{Division and Revision of Responsibilities}
\vspace{-5pt}
\begin{itemize}
    \setlength{\itemsep}{0pt}
    \setlength{\parskip}{0pt}
    \setlength{\parsep}{0pt}
    \item \textbf{Ajay:} Data cleaning operations, summary tables, distribution plots (including histograms), and outlier detection checks.
    \item \textbf{Rushmith:} Distribution fitting, parameter estimation, constructing confidence intervals, and running permutation simulations.
    \item \textbf{Bhargav:} Setting up regression models, conducting diagnostic checks, and analyzing screen-count mediation.
    \item \textbf{Aryan:} Formal hypothesis testing, calculating effect sizes, multiple-testing corrections, and revenue concentration analysis.
\end{itemize}
\vspace{-6pt}

\textbf{Updated Tracking Spreadsheet}\\
\href{https://docs.google.com/spreadsheets/d/1Uq2D5MdjAsLFLXG8qR6Nv4SFCqRq0ZuyhnNrdCJzuk0/edit?usp=sharing}{Google Sheets Project Tracker}\\
We have updated our tracker to include a dedicated row containing our meta-information content as per the guidelines. Edit
access has been fully verified and granted to our mentor, Avadh Hindocha, ensuring seamless collaboration.

\textbf{Blockers and Challenges}\\
Our primary challenge is the absence of exact release years in the dataset, which prevents us from adjusting revenues for
inflation and identifying strict temporal trends. Additionally, the histograms revealed severe right-skewness and outliers
in the revenue columns. We are consulting statistical resources to decide whether these extreme values should be capped,
log-transformed, or analyzed separately to prevent them from overly skewing our regression models. Immediate guidance
from the Teaching Team on handling inflation without precise release years would be highly appreciated

\end{document}
